\subsection{trumpet， חֲצֹצְרֹת，שׁוֺפָר}
חֲצֹצְרָה noun feminine clarion (Late Hebrew חֲצוֺצֶרֶת, Aramaic חֲצוֺצַרְתָּצ) — mostly P and late; — absolute ׳ח Hosea 5:8; plural absolute חֲצֹצְרוֺת Numbers 10:8 22t.; חֲצֹצְרֹת Numbers 10:9,10; construct id. Numbers 31:6; 2Chronicles 13:12; חֲצוֺצְרֹת Numbers 10:2; clarion


as secular instr. Hosea 5:8 ("" שׁוֺפָר) 2 Kings 11:14 (twice in verse) = 2Chronicles 23:13 (twice in verse).

2 as sacred instr. 2 Kings 12:14, especially for use by priests (only P, Psalm 98 and Chronicles).

a. ׳תקע בח (of blowing a single long blast) Numbers 10:3,4,7,8, to gather congregation or ׳נשׂיא together, and, on festivals, over sacrifice, 'to be remembered before ׳י,' Numbers 10:10.

b. ׳תקע תרועה בח, or ׳הריע בח (of sounding alarm, — a series of quick blasts) for camps to move Numbers 10:5,6, also in battle, Numbers 10:9, 'to be remembered before ׳י;' so Numbers 31:6; 2Chronicles 13:12 (compare 2 Chronicles 13:14), both התרועה ׳ח.

c. especially in Chr's descriptions of ceremonies at festivals, to express rejoicing: 1 Chronicles 13:8; 1 Chronicles 15:28 ("" קוֺל שׁוֺפָר), 1 Chronicles 16:6,42; 2Chronicles 15:14 ("" שׁוֺפָר), 2 Chronicles 20:28; 29:26,27; Ezra 3:10; Nehemiah 12:35,41; Psalm 98:6 ("" קוֺל שׁוֺפָר); ׳הָרִים קוֺל בח2Chronicles 5:13; ׳מחצצרים בח 1 Chronicles 15:24; 2Chronicles 5:12,13; 13:14; in 2 Chronicles 29:28 this participle agrees with noun in sense, and is masculine; and the clarions (= players on the clarions) sounded. — The חֲצֹצְרָה, or (sacred) clarion, was a long, straight, slender metal tube, with flaring end, see BenzArchäol. 277; distinguished thus from the שׁוֺפָר which was originally a ram's horn, and probably always retained the horn-shape; the שׁוֺפָר is mentioned constantly in the earlier literature, and was used by watchmen, warriors, etc., as well as priests (see Benzib. 276 and שׁוֺפָר).

耶和華曉諭摩西說：「你要用銀子做兩枝號，都要錘出來的，用以招聚會眾，並叫眾營起行。吹這號的時候，全會眾要到你那裏，聚集在會幕門口。若單吹一枝，眾首領，就是以色列軍中的統領，要聚集到你那裏。吹出大聲的時候，東邊安的營都要起行。二次吹出大聲的時候，南邊安的營都要起行。他們將起行，必吹出大聲。但招聚會眾的時候，你們要吹號，卻不要吹出大聲。亞倫子孫作祭司的要吹這號；這要作你們世世代代永遠的定例。你們在自己的地，與欺壓你們的敵人打仗，就要用號吹出大聲，便在耶和華──你們的神面前得蒙紀念，也蒙拯救脫離仇敵。在你們快樂的日子和節期，並月朔，獻燔祭和平安祭，也要吹號，這都要在你們的神面前作為紀念，我是耶和華──你們的神。」（民數記十1～10）


另外，在特定的日子、節期和月朔，歡樂的場合中，他們也是吹奏銀號助興；就是每逢快樂的場合、各種節期、每月初一、每年的吹角節（民數記廿九1；利未記廿三24）、禧年的七月初十（利未記廿五9），各節期的宗教活動，也都要吹號。


吹號乃傳達最高神的命令，吹號的意義即表明傳揚神的聖言（參：馬可福音十六15；啟示錄十四6、7）。端視今日的基督徒在屬靈的曠野途中，有作戰，也有歡樂，有前進，也有等待，我們也是要藉著基督的號令，從四方招聚祂的百姓，並藉著基督的聖言，來帶領教會和幫助選民，祂是銀號的真體，是真理的發言者。

然而為宣告救贖之道，耶穌基督就像銀號一般，經由千錘百鍊，祂受盡了千辛萬苦，被神擊打苦待，才得以成全（希伯來書九12；以賽亞書五十三4～10），其目的就在宣揚救人之道，並對世人傳警告（參：馬太福音廿28；使徒行傳十三26）。


在歷代之中，選民所備的號並沒有局限幾枝；摩西之曠野時代，是選民最先設立吹號制度，他依耶和華神的吩咐，預備的銀號是兩枝。到了大衛的統一王國時期，就備有七枝號筒了（參：歷代志上十五24）。及至以色列國全盛時期，當時的所羅門王一口氣製造了120枝的號筒（歷代志下五12）。甚至末後基督再臨時，祂要差遣使者，用號筒的大聲，將祂的選民，從四方，從天這邊到天那邊，都招聚了來（參：馬太福音廿四31）。可見隨著時代的改變，不但使用號筒的數量跟著增加，而且也成了基督再臨的宣告。

因此，基督是那歷史以來，真正的呼召者，祂在古時以銀號呼召祂的百姓，今日更親自藉著教會，呼召末後的選民，祂的救贖的號聲是永不停頓，直到末日祂的再臨（啟示錄七9～17；馬太福音廿四30～31）。
今日的教會，必須嚴守歸回聖經的原則，傾聽基督的真理，分辨主的號聲，以恢復使徒時代的教會為標的，接受全備的教理，領受上頭澆灌的聖靈，確實傳揚神的聖言（馬可福音十六15；啟示錄十四6～7），

以西結書卅三1～9）。因此，要積極宣揚救贖之訊息，努力鼓吹，高舉基督的號角，盡祭司之職責（哥林多前書九16、17、19～22）











\subsection{愤怒 אַ֨ף  nostril, nose, face, anger（oftener divine）}
\subsubsection{其他书卷 אַ֨ף  nostril, nose, face, anger（oftener divine）}
\begin{itemize} 
\itemsep-.25em
\item 3 mostly anger, human Genesis 27:45; Genesis 49:6,7 + (45 t.); oftener divine Exodus 32:12; Deuteronomy 9:19; 2 Kings 24:20 + (177 t.); often subject חָרָה (וַיִּ֫חַר etc.) his anger was kindled Genesis 30:2; Genesis 39:19; Exodus 4:14; Exodus 22:23; Exodus 32:10,11 +; in various combinations, especially חֲרוֺן אַף fierceness of anger Exodus 32:12; Numbers 25:4; Numbers 32:14 +; compare חֳרִיאָֿף 1 Samuel 20:34; בַּעַלאַֿף Proverbs 22:24 one given to anger, etc.; אֶרֶךְ אַמַּיִם slow to anger Exodus 34:6; Numbers 14:18; Nehemiah 9:17 7t. of God; Proverbs 14:29; Proverbs 15:18; Proverbs 16:32; Proverbs 25:15 of man.
\item 出埃及記 22:24 並要發烈怒，用刀殺你們，使你們的妻子為寡婦，兒女為孤兒。
\item 民數記 16:46 摩西對亞倫說：「拿你的香爐，把壇上的火盛在其中，又加上香，快快帶到會眾那裡，為他們贖罪。因為有憤怒從耶和華那裡出來，瘟疫已經發作了。」
民數記 18:5 你們要看守聖所和壇，免得憤怒再臨到以色列人。
\item 歷代志下 36:16 他們卻戲笑神的使者，藐視他的言語，譏誚他的先知，以致耶和華的憤怒向他的百姓發作，無法可救。
\end{itemize}  

\subsubsection{詩篇}
\begin{itemize} 
\itemsep-.25em
\item 詩篇 88:7 你的憤怒重壓我身，你用一切的波浪困住我。（細拉）
\item 詩篇 89:46 耶和華啊，這要到幾時呢？你要將自己隱藏到永遠嗎？你的憤怒如火焚燒要到幾時呢？箴言 16:14 王的震怒如殺人的使者，但智慧人能止息王怒。
\item 詩篇 106:40 所以耶和華的怒氣向他的百姓發作，憎惡他的產業，
\item 詩篇 124:3 向我們發怒的時候，就把我們活活地吞了。
\end{itemize}  

\subsubsection{以賽亞書，耶利米哀歌，哈巴谷書}
\begin{itemize} 
\itemsep-.25em
\item 以賽亞書 13:13 我萬軍之耶和華在憤恨中發烈怒的日子，必使天震動，使地搖撼離其本位。
\item 以賽亞書 60:10 「外邦人必建築你的城牆，他們的王必服侍你。我曾發怒擊打你，現今卻施恩憐恤你。
\item 耶利米哀歌 3:1 我是因耶和華憤怒的杖遭遇困苦的人。
\item agitation, excitement, raging; — ׳ר absolute Habakkuk 3:2 +, construct Job 37:2; suffix רָנְזֶךָ֑ Isaiah 14:3; — raging Job 3:17; disquiet, turmoil Isaiah 14:3; Job 3:26; Job 14:1; raging, wrath Habakkuk 3:2; קֹלוֺ ׳ר Job 37:2 rumbling of his voice (i.e. thunder); of excitement of warhorse, בְּרַעַשׁ וְרֹגֶּז Job 39:24.
哈巴谷書 3:2 耶和華啊，我聽見你的名聲就懼怕。耶和華啊，求你在這些年間復興你的作為，在這些年間顯明出來，在發怒 的時候以憐憫為念。
\end{itemize}  


\subsubsection{羅馬書 2:5，9:22}
\begin{itemize} 
\itemsep-.25em
\item 羅馬書 2:5 你竟任著你剛硬不悔改的心，為自己積蓄憤怒，以致神震怒，顯他公義審判的日子來到。
\item 窯匠難道沒有權柄從一團泥裡拿一塊做成貴重的器皿，又拿一塊做成卑賤的器皿嗎？ 22倘若神要顯明他的憤怒，彰顯他的權能，就多多忍耐寬容那可怒、預備遭毀滅的器皿， 23又要將他豐盛的榮耀彰顯在那蒙憐憫、早預備得榮耀的器皿上
\end{itemize}  

\subsection{患难，大患难2347: $\theta$$\lambda$$\widetilde{\iota}$$\psi$$\iota$$\varsigma$ θλῖψις}


\subsubsection{馬太福音 13:17-23 $\theta$$\lambda$$\acute{\iota}$$\psi$$\epsilon$$\omega$$\zeta$,24:3-14 $\theta$$\lambda$$\tilde{\iota}$$\psi$$\iota$$\nu$,24:20-31 $\theta$$\lambda$$\tilde{\iota}$$\psi$$\iota$$\zeta$ $\theta$$\lambda$$\tilde{\iota}$$\psi$$\iota$$\nu$}    
\begin{itemize} 
\itemsep-.25em
\item 17 「我實在告訴你們：從前有許多先知和義人要看你們所看的，卻沒有看見；要聽你們所聽的，卻沒有聽見。 18 所以，你們當聽這撒種的比喻： 19 凡聽見天國道理不明白的，那惡者就來，把所撒在他心裡的奪了去，這就是撒在路旁的了。 20 撒在石頭地上的，就是人聽了道，當下歡喜領受， 21 只因心裡沒有根，不過是暫時的，及至為道遭了\underline{患難$\theta$$\lambda$$\acute{\iota}$$\psi$$\epsilon$$\omega$$\zeta$}，或是受了逼迫，立刻就跌倒了。 22 撒在荊棘裡的，就是人聽了道，後來有世上的思慮、錢財的迷惑把道擠住了，不能結實。 23 撒在好地上的，就是人聽道明白了，後來結實，有一百倍的，有六十倍的，有三十倍的。」   
\item 耶穌在橄欖山上坐著，門徒暗暗地來，說：「請告訴我們，什麼時候有這些事？你降臨和世界的末了有什麼預兆呢？」 4 耶穌回答說：「你們要謹慎，免得有人迷惑你們。 5 因為將來有好些人冒我的名來，說：『我是基督』，並且要迷惑許多人。 6 你們也要聽見打仗和打仗的風聲，總不要驚慌，因為這些事是必須有的，只是末期還沒有到。 7 民要攻打民，國要攻打國，多處必有饑荒、地震， 8 這都是災難的起頭。 9 那時，人要把你們陷在\underline{患難$\theta$$\lambda$$\tilde{\iota}$$\psi$$\iota$$\nu$}裡，也要殺害你們；你們又要為我的名被萬民恨惡。「那時，必有許多人跌倒，也要彼此陷害、彼此恨惡， 11 且有好些假先知起來，迷惑多人。 12 只因不法的事增多，許多人的愛心才漸漸冷淡了。 13 唯有忍耐到底的，必然得救。 14 這天國的福音要傳遍天下，對萬民作見證，然後末期才來到。
\item 20 「你們應當祈求，叫你們逃走的時候，不遇見冬天或是安息日。 21 因為那時必有\underline{大災難$\theta$$\lambda$$\tilde{\iota}$$\psi$$\iota$$\zeta$}，從世界的起頭直到如今沒有這樣的災難，後來也必沒有。 22 若不減少那日子，凡有血氣的總沒有一個得救的；只是為選民，那日子必減少了。 23 那時，若有人對你們說『基督在這裡』，或說『基督在那裡』，你們不要信。 24 因為假基督、假先知將要起來，顯大神蹟、大奇事，倘若能行，連選民也就迷惑了。 25 看哪，我預先告訴你們了！ 26 若有人對你們說『看哪，基督在曠野裡』，你們不要出去；或說『看哪，基督在內屋中』，你們不要信。 27 閃電從東邊發出，直照到西邊，人子降臨也要這樣。 28 屍首在哪裡，鷹也必聚在那裡。「那些日子的\underline{災難$\theta$$\lambda$$\tilde{\iota}$$\psi$$\iota$$\nu$}一過去，『日頭就變黑了，月亮也不放光，眾星要從天上墜落，天勢都要震動。』 30 那時，人子的兆頭要顯在天上，地上的萬族都要哀哭。他們要看見人子有能力、有大榮耀駕著天上的雲降臨。 31 他要差遣使者，用號筒的大聲，將他的選民從四方，從天這邊到天那邊，都招聚了來。
\end{itemize}   
    
 \subsubsection{馬可福音 4:13-20 $\theta$$\lambda$$\acute{\iota}$$\psi$$\epsilon$$\omega$$\zeta$,13:14-27 $\theta$$\lambda$$\tilde{\iota}$$\psi$$\iota$$\zeta$, $\theta$$\lambda$$\tilde{\iota}$$\psi$$\iota$$\nu$} 
 \begin{itemize} 
\itemsep-.25em
\item {馬可福音 4:13-20}\\
13 又對他們說：「你們不明白這比喻嗎？這樣怎能明白一切的比喻呢？ 14 撒種之人所撒的就是道。 15 那撒在路旁的，就是人聽了道，撒旦立刻來，把撒在他心裡的道奪了去。 16 那撒在石頭地上的，就是人聽了道，立刻歡喜領受， 17 但他心裡沒有根，不過是暫時的，及至為道遭了\underline{患難$\theta$$\lambda$$\acute{\iota}$$\psi$$\epsilon$$\omega$$\zeta$}，或是受了逼迫，立刻就跌倒了。 18 還有那撒在荊棘裡的，就是人聽了道，19 後來有世上的思慮、錢財的迷惑和別樣的私慾進來，把道擠住了，就不能結實。 20 那撒在好地上的，就是人聽道又領受，並且結實，有三十倍的，有六十倍的，有一百倍的。」
\item{馬可福音 13:14-27}   \\
「你們看見『那行毀壞可憎的』站在不當站的地方——讀這經的人須要會意——那時，在猶太的，應當逃到山上； 15 在房上的，不要下來，也不要進去拿家裡的東西； 16 在田裡的，也不要回去取衣裳。 17 當那些日子，懷孕的和奶孩子的有禍了！ 18 你們應當祈求，叫這些事不在冬天臨到。 19 因為在那些日子必有\underline{災難$\theta$$\lambda$$\tilde{\iota}$$\psi$$\iota$$\zeta$}，自從神創造萬物直到如今並沒有這樣的災難，後來也必沒有。 20 若不是主減少那日子，凡有血氣的總沒有一個得救的；只是為主的選民，他將那日子減少了。 21 那時，若有人對你們說『看哪，基督在這裡』，或說『基督在那裡』，你們不要信。 22 因為假基督、假先知將要起來，顯神蹟奇事，倘若能行，就把選民迷惑了。 23 你們要謹慎！看哪，凡事我都預先告訴你們了！24 「在那些日子，那\underline{災難$\theta$$\lambda$$\tilde{\iota}$$\psi$$\iota$$\nu$}以後，『日頭要變黑了，月亮也不放光， 25 眾星要從天上墜落，天勢都要震動。』 26 那時，他們 b要看見人子有大能力、大榮耀駕雲降臨。 27 他要差遣天使，把他的選民從四方，從地極直到天邊，都招聚了來。  
\end{itemize}  

\subsubsection{約翰福音 16：19-33 $\theta$$\lambda$$\acute{\iota}$$\psi$$\epsilon$$\omega$$\varsigma$ $\theta$$\lambda$$\tilde{\iota}$$\psi$$\iota$$\nu$}
耶穌看出他們要問他，就說：「我說『等不多時，你們就不得見我；再等不多時，你們還要見我』，你們為這話彼此相問嗎？ 20 我實實在在地告訴你們：你們將要痛哭、哀號，世人倒要喜樂。你們將要憂愁，然而你們的憂愁要變為喜樂。 21 婦人生產的時候就憂愁，因為她的時候到了；既生了孩子，就不再記念那\underline{苦楚$\theta$$\lambda$$\acute{\iota}$$\psi$$\epsilon$$\omega$$\varsigma$}，因為歡喜世上生了一個人。 22 你們現在也是憂愁，但我要再見你們，你們的心就喜樂了，這喜樂也沒有人能奪去。23 「到那日，你們什麼也就不問我了。我實實在在地告訴你們：你們若向父求什麼，他必因我的名賜給你們。 24 向來你們沒有奉我的名求什麼，如今你們求，就必得著，叫你們的喜樂可以滿足。25 「這些事，我是用比喻對你們說的。時候將到，我不再用比喻對你們說，乃要將父明明地告訴你們。 26 到那日，你們要奉我的名祈求。我並不對你們說，我要為你們求父； 27 父自己愛你們，因為你們已經愛我，又信我是從父出來的。 28 我從父出來，到了世界；我又離開世界，往父那裡去。」 29 門徒說：「如今你是明說，並不用比喻了。 30 現在我們曉得你凡事都知道，也不用人問你，因此我們信你是從神出來的。」 31 耶穌說：「現在你們信嗎？ 32 看哪，時候將到，且是已經到了，你們要分散，各歸自己的地方去，留下我獨自一人；其實我不是獨自一人，因為有父與我同在。 33 我將這些事告訴你們，是要叫你們在我裡面有平安。在世上你們有\underline{苦難$\theta$$\lambda$$\tilde{\iota}$$\psi$$\iota$$\nu$}，但你們可以放心，我已經勝了世界。」


\subsubsection{使徒行傳 7:9-12 $\theta$$\lambda$$\tilde{\iota}$$\psi$$\iota$$\zeta$, 14:19-23 $\theta$$\lambda$$\acute{\iota}$$\psi$$\epsilon$$\omega$$\nu$, 20:17-24 $\theta$$\lambda$$\acute{\iota}$$\psi$$\epsilon$$\iota$$\zeta$} 
\begin{itemize} 
\itemsep-.25em
\item{使徒行傳 7:9-12}\\  9 先祖嫉妒約瑟，把他賣到埃及去。神卻與他同在， 10 救他脫離一切苦難，又使他在埃及王法老面前得恩典、有智慧。法老就派他做埃及國的宰相兼管全家。 11 後來埃及和迦南全地遭遇饑荒，大受\underline{艱難$\theta$$\lambda$$\tilde{\iota}$$\psi$$\iota$$\zeta$}，我們的祖宗就絕了糧。 12 雅各聽見在埃及有糧，就打發我們的祖宗初次往那裡去。
\item{使徒行傳 14:19-23} \\
 但有些猶太人從安提阿和以哥念來，挑唆眾人，就用石頭打保羅，以為他是死了，便拖到城外。 20 門徒正圍著他，他就起來，走進城去。第二天，同巴拿巴往特庇去。 21 對那城裡的人傳了福音，使好些人做門徒，就回路司得、以哥念、安提阿去， 22 堅固門徒的心，勸他們恆守所信的道，又說：「我們進入神的國，必須經歷許多\underline{艱難$\theta$$\lambda$$\acute{\iota}$$\psi$$\epsilon$$\omega$$\nu$}。」 23 二人在各教會中選立了長老，又禁食禱告，就把他們交託所信的主。
\item{使徒行傳 20:17-24} \\
保羅從米利都打發人往以弗所去，請教會的長老來。 18 他們來了，保羅就說：「你們知道自從我到亞細亞的日子以來，在你們中間始終為人如何， 19 服侍主凡事謙卑，眼中流淚，又因猶太人的謀害經歷試煉。 20 你們也知道，凡於你們有益的，我沒有一樣避諱不說的，或在眾人面前，或在各人家裡，我都教導你們。 21 又對猶太人和希臘人證明當向神悔改，信靠我主耶穌基督。 22 現在我往耶路撒冷去，心甚迫切，不知道在那裡要遇見什麼事。 23 但知道聖靈在各城裡向我指證，說有捆鎖與\underline{患難$\theta$$\lambda$$\acute{\iota}$$\psi$$\epsilon$$\iota$$\zeta$}等待我。 24 我卻不以性命為念，也不看為寶貴，只要行完我的路程，成就我從主耶穌所領受的職事，證明神恩惠的福音。 
 \end{itemize}
 
 
 \subsubsection{羅馬書 2：1-11 $\theta$$\lambda$$\widetilde{\iota}$$\psi$$\iota$$\varsigma$, 5:1-5 $\theta$$\lambda$$\acute{\iota}$$\psi$$\epsilon$$\sigma$$\iota$$\nu$ $\theta$$\lambda$$\tilde{\iota}$$\psi$$\iota$$\zeta$, 8:31-39 $\theta$$\lambda$$\tilde{\iota}$$\psi$$\iota$$\zeta$, 12:9-16 $\theta$$\lambda$$\acute{\iota}$$\psi$$\epsilon$$\iota$}
 \begin{itemize} 
\itemsep-.25em
\item{羅馬書 2：1-11}\\1 你這論斷人的，無論你是誰，也無可推諉。你在什麼事上論斷人，就在什麼事上定自己的罪，因你這論斷人的，自己所行卻和別人一樣。 2 我們知道這樣行的人，神必照真理審判他。 3 你這人哪，你論斷行這樣事的人，自己所行的卻和別人一樣，你以為能逃脫神的審判嗎？ 4 還是你藐視他豐富的恩慈、寬容、忍耐，不曉得他的恩慈是領你悔改呢？ 5 你竟任著你剛硬不悔改的心，為自己積蓄憤怒，以致神震怒，顯他公義審判的日子來到。 6 他必照各人的行為報應各人。 7 凡恆心行善，尋求榮耀、尊貴和不能朽壞之福的，就以永生報應他們； 8 唯有結黨、不順從真理反順從不義的，就以憤怒、惱恨報應他們。 9 將\underline{患難$\theta$$\lambda$$\widetilde{\iota}$$\psi$$\iota$$\varsigma$}、困苦加給一切作惡的人，先是猶太人，後是希臘人； 10 卻將榮耀、尊貴、平安加給一切行善的人，先是猶太人，後是希臘人。 11 因為神不偏待人。
\item{羅馬書 5:1-5} \\我們既因信稱義，就藉著我們的主耶穌基督得與神相和。 2 我們又藉著他，因信得進入現在所站的這恩典中，並且歡歡喜喜盼望神的榮耀。 3 不但如此，就是在\underline{患難$\theta$$\lambda$$\acute{\iota}$$\psi$$\epsilon$$\sigma$$\iota$$\nu$}中也是歡歡喜喜的。因為知道\underline{患難$\theta$$\lambda$$\tilde{\iota}$$\psi$$\iota$$\zeta$}生忍耐， 4 忍耐生老練，老練生盼望， 5 盼望不至於羞恥，因為所賜給我們的聖靈將神的愛澆灌在我們心裡。  
\item{羅馬書 8:31-39}   \\ 31 既是這樣，還有什麼說的呢？神若幫助我們，誰能敵擋我們呢？ 32 神既不愛惜自己的兒子，為我們眾人捨了，豈不也把萬物和他一同白白地賜給我們嗎？ 33 誰能控告神所揀選的人呢？有神稱他們為義了。 c 34 誰能定他們的罪呢？有基督耶穌已經死了，而且從死裡復活，現今在神的右邊，也替我們祈求。35 誰能使我們與基督的愛隔絕呢？難道是\underline{患難$\theta$$\lambda$$\tilde{\iota}$$\psi$$\iota$$\zeta$}嗎？是困苦嗎？是逼迫嗎？是飢餓嗎？是赤身露體嗎？是危險嗎？是刀劍嗎？ 36 如經上所記：「我們為你的緣故終日被殺，人看我們如將宰的羊。」 37 然而，靠著愛我們的主，在這一切的事上已經得勝有餘了。 38 因為我深信：無論是死，是生，是天使，是掌權的，是有能的，是現在的事，是將來的事， 39 是高處的，是低處的，是別的受造之物，都不能叫我們與神的愛隔絕；這愛是在我們的主基督耶穌裡的。 
 \item{羅馬書 12:9-16} \\ 愛人不可虛假；惡要厭惡，善要親近。 10 愛弟兄，要彼此親熱；恭敬人，要彼此推讓。 11 殷勤不可懶惰；要心裡火熱，常常服侍主。 12 在指望中要喜樂，在\underline{患難$\theta$$\lambda$$\acute{\iota}$$\psi$$\epsilon$$\iota$}中要忍耐，禱告要恆切。 13 聖徒缺乏要幫補，客要一味地款待。 14 逼迫你們的，要給他們祝福；只要祝福，不可咒詛。 15 與喜樂的人要同樂，與哀哭的人要同哭。 16 要彼此同心；不要志氣高大，倒要俯就卑微的人；不要自以為聰明。 
\end{itemize}
 
 
\subsubsection{哥林多前書 7：26-31 $\theta$$\lambda$$\tilde{\iota}$$\psi$$\iota$$\nu$} 
26 因現今的艱難，據我看來，人不如守素安常才好。 27 你有妻子纏著呢，就不要求脫離；你沒有妻子纏著呢，就不要求妻子。 28 你若娶妻，並不是犯罪；處女若出嫁，也不是犯罪。然而這等人肉身必受\underline{苦難$\theta$$\lambda$$\tilde{\iota}$$\psi$$\iota$$\nu$}，我卻願意你們免這苦難。 29 弟兄們，我對你們說，時候減少了。從此以後，那有妻子的，要像沒有妻子； 30 哀哭的，要像不哀哭；快樂的，要像不快樂；置買的，要像無有所得； 31 用世物的，要像不用世物；因為這世界的樣子將要過去了。
 
 
\subsubsection{林後 1:3-11 $\theta$$\lambda$$\acute{\iota}$$\psi$$\epsilon$$\iota$ $\theta$$\lambda$$\acute{\iota}$$\psi$$\epsilon$$\omega$$\zeta$,  2:1-4/4:16-18 $\theta$$\lambda$$\acute{\iota}$$\psi$$\epsilon$$\omega$$\zeta$, 6:1-10 $\theta$$\lambda$$\acute{\iota}$$\psi$$\epsilon$$\sigma$$\iota$$\nu$,  7:2-4 $\theta$$\lambda$$\acute{\iota}$$\psi$$\epsilon$$\iota$, 8:1-8 $\theta$$\lambda$$\acute{\iota}$$\psi$$\epsilon$$\omega$$\zeta$}
\begin{itemize} 
\itemsep-.25em
\item{哥林多後書 1：3-11}\\ 願頌讚歸於我們的主耶穌基督的父神，就是發慈悲的父，賜各樣安慰的神！ 4 我們在一切\underline{患難$\theta$$\lambda$$\acute{\iota}$$\psi$$\epsilon$$\iota$}中，他就安慰我們，叫我們能用神所賜的安慰去安慰那遭各樣患難的人。 5 我們既多受基督的苦楚，就靠基督多得安慰。 6 我們受患難呢，是為叫你們得安慰、得拯救；我們得安慰呢，也是為叫你們得安慰。這安慰能叫你們忍受我們所受的那樣苦楚。 7 我們為你們所存的盼望是確定的，因為知道你們既是同受苦楚，也必同得安慰。 8 弟兄們，我們不要你們不曉得，我們從前在亞細亞遭遇\underline{苦難$\theta$$\lambda$$\acute{\iota}$$\psi$$\epsilon$$\omega$$\zeta$}，被壓太重，力不能勝，甚至連活命的指望都絕了， 9 自己心裡也斷定是必死的，叫我們不靠自己，只靠叫死人復活的神。 10 他曾救我們脫離那極大的死亡，現在仍要救我們，並且我們指望他將來還要救我們。 11 你們以祈禱幫助我們，好叫許多人為我們謝恩，就是為我們因許多人所得的恩。
\item{哥林多後書 2：1-4} \\ 1 我自己定了主意，再到你們那裡去，必須大家沒有憂愁。 2 倘若我叫你們憂愁，除了我叫那憂愁的人以外，誰能叫我快樂呢？ 3 我曾把這事寫給你們，恐怕我到的時候，應該叫我快樂的那些人反倒叫我憂愁。我也深信，你們眾人都以我的快樂為自己的快樂。 4 我先前心裡\underline{難過$\theta$$\lambda$$\acute{\iota}$$\psi$$\epsilon$$\omega$$\zeta$}痛苦，多多地流淚寫信給你們，不是叫你們憂愁，乃是叫你們知道我格外地疼愛你們。
\item{哥林多後書 4：16-18} \\ 16 所以我們不喪膽，外體雖然毀壞，內心卻一天新似一天。 17 我們這至暫至輕的\underline{苦楚$\theta$$\lambda$$\acute{\iota}$$\psi$$\epsilon$$\omega$$\zeta$}，要為我們成就極重無比、永遠的榮耀。 18 原來我們不是顧念所見的，乃是顧念所不見的，因為所見的是暫時的，所不見的是永遠的。     
\item{哥林多後書 6：1-10} \\1 我們與神同工的，也勸你們不可徒受他的恩典。 2 因為他說：「在悅納的時候，我應允了你；在拯救的日子，我搭救了你。」看哪，現在正是悅納的時候！現在正是拯救的日子！ 3 我們凡事都不叫人有妨礙，免得這職分被人毀謗； 4 反倒在各樣的事上表明自己是神的用人，就如在許多的忍耐，\underline{患難$\theta$$\lambda$$\acute{\iota}$$\psi$$\epsilon$$\sigma$$\iota$$\nu$}，窮乏，困苦， 5 鞭打，監禁，擾亂，勤勞，警醒，不食， 6 廉潔，知識，恆忍，恩慈，聖靈的感化，無偽的愛心， 7 真實的道理，神的大能；仁義的兵器在左在右， 8 榮耀羞辱，惡名美名；似乎是誘惑人的，卻是誠實的； 9 似乎不為人所知，卻是人所共知的；似乎要死，卻是活著的；似乎受責罰，卻是不致喪命的； 10 似乎憂愁，卻是常常快樂的；似乎貧窮，卻是叫許多人富足的；似乎一無所有，卻是樣樣都有的。    
 \item{哥林多後書 7：2-4}   \\ 2 你們要心地寬大收納我們。我們未曾虧負誰，未曾敗壞誰，未曾占誰的便宜。 3 我說這話不是要定你們的罪，我已經說過，你們常在我們心裡，情願與你們同生同死。 4 我大大地放膽，向你們說話；我因你們多多誇口，滿得安慰，我們在一切\underline{患難$\theta$$\lambda$$\acute{\iota}$$\psi$$\epsilon$$\iota$}中分外地快樂。 
 \item{哥林多後書 8：1-8}   \\弟兄們，我把神賜給馬其頓眾教會的恩告訴你們， 2 就是他們在患難中受\underline{大試煉$\theta$$\lambda$$\acute{\iota}$$\psi$$\epsilon$$\omega$$\zeta$}的時候，仍有滿足的快樂，在極窮之間還格外顯出他們樂捐的厚恩。 3 我可以證明他們是按著力量，而且也過了力量，自己甘心樂意地捐助， 4 再三地求我們，准他們在這供給聖徒的恩情上有份。 5 並且他們所做的，不但照我們所想望的，更照神的旨意，先把自己獻給主，又歸附了我們。 6 因此我勸提多，既然在你們中間開辦這慈惠的事，就當辦成了。 7 你們既然在信心、口才、知識、熱心和待我們的愛心上，都格外顯出滿足來，就當在這慈惠的事上也格外顯出滿足來。 8 我說這話，不是吩咐你們，乃是藉著別人的熱心試驗你們愛心的實在。    
\end{itemize}
 
 
 
 \subsubsection{以弗所書 3:2-13 $\theta$$\lambda$$\acute{\iota}$$\psi$$\epsilon$$\sigma$$\acute{\iota}$$\nu$} 
 2 諒必你們曾聽見神賜恩給我，將關切你們的職分託付我， 3 用啟示使我知道福音的奧祕，正如我以前略略寫過的。 4 你們念了，就能曉得我深知基督的奧祕。 5 這奧祕在以前的世代沒有叫人知道，像如今藉著聖靈啟示他的聖使徒和先知一樣。 6 這奧祕就是外邦人在基督耶穌裡，藉著福音，得以同為後嗣，同為一體，同蒙應許。 7 我做了這福音的執事，是照神的恩賜，這恩賜是照他運行的大能賜給我的。 8 我本來比眾聖徒中最小的還小，然而他還賜我這恩典，叫我把基督那測不透的豐富傳給外邦人， 9 又使眾人都明白，這歷代以來隱藏在創造萬物之神裡的奧祕是如何安排的， 10 為要藉著教會，使天上執政的、掌權的現在得知神百般的智慧。 11 這是照神從萬世以前在我們主基督耶穌裡所定的旨意。 12 我們因信耶穌，就在他裡面放膽無懼，篤信不疑地來到神面前。 13 所以，我求你們不要因我為你們所受的\underline{患難$\theta$$\lambda$$\acute{\iota}$$\psi$$\epsilon$$\sigma$$\acute{\iota}$$\nu$}喪膽，這原是你們的榮耀。
 
 
 \subsubsection{腓立比書 1：12-17 $\theta$$\lambda$$\widetilde{\iota}$$\psi$$\iota$$\nu$, 4：10-14 $\theta$$\lambda$$\acute{\iota}$$\psi$$\epsilon$$\iota$}
 \begin{itemize} 
\itemsep-.25em
\item{腓立比書 1：12-17}\\ 12 弟兄們，我願意你們知道，我所遭遇的事更是叫福音興旺， 13 以致我受的捆鎖在御營全軍和其餘的人中，已經顯明是為基督的緣故； 14 並且那在主裡的弟兄，多半因我受的捆鎖就篤信不疑，越發放膽傳神的道，無所懼怕。 15 有的傳基督是出於嫉妒紛爭，也有的是出於好意。 16 這一等是出於愛心，知道我是為辯明福音設立的； 17 那一等傳基督是出於結黨，並不誠實，意思要加增我捆鎖的\underline{苦楚$\theta$$\lambda$$\widetilde{\iota}$$\psi$$\iota$$\nu$}。
 \begin{cjhebrew}
hr.s  .S.hl  r.s  צַר, לַחַץ,צָרָה 
 \end{cjhebrew}
 \item{腓立比書 4：10-14}\\ 10 我靠主大大地喜樂，因為你們思念我的心如今又發生；你們向來就思念我，只是沒得機會。 11 我並不是因缺乏說這話，我無論在什麼景況都可以知足，這是我已經學會了。 12 我知道怎樣處卑賤，也知道怎樣處豐富，或飽足或飢餓，或有餘或缺乏，隨事隨在，我都得了祕訣。 13 我靠著那加給我力量的，凡事都能做。 14 然而，你們和我同受\underline{患難$\theta$$\lambda$$\acute{\iota}$$\psi$$\epsilon$$\iota$}原是美事。
 \end{itemize}
 
 
 \subsubsection{歌羅西書 1：23-27 $\theta$$\lambda$$\acute{\iota}$$\psi$$\epsilon$$\omega$$\nu$}
 23 只要你們在所信的道上恆心，根基穩固，堅定不移，不致被引動失去福音的盼望。這福音就是你們所聽過的，也是傳於普天下萬人聽的。我保羅也做了這福音的執事。24 現在我為你們受苦，倒覺歡樂，並且為基督的身體，就是為教會，要在我肉身上補滿基督\underline{患難$\theta$$\lambda$$\acute{\iota}$$\psi$$\epsilon$$\omega$$\nu$}的缺欠。 25 我照神為你們所賜我的職分做了教會的執事，要把神的道理傳得全備。 26 這道理就是歷世歷代所隱藏的奧祕，但如今向他的聖徒顯明了。 27 神願意叫他們知道，這奧祕在外邦人中有何等豐盛的榮耀，就是基督在你們心裡成了有榮耀的盼望。
 
 \subsubsection{帖前 1:2-8 $\theta$$\lambda$$\acute{\iota}$$\psi$$\epsilon$$\iota$, 3:1-10 $\theta$$\lambda$$\acute{\iota}$$\psi$$\epsilon$$\sigma$$\iota$$\nu$ $\theta$$\lambda$$\acute{\iota}$$\beta$$\epsilon$$\sigma$$\theta$$\alpha$$\iota$$\tau$$o$ $\theta$$\lambda$$\acute{\iota}$$\psi$$\epsilon$$\iota$, 帖後 1:3-10 $\theta$$\lambda$$\acute{\iota}$$\psi$$\epsilon$$\sigma$$\iota$$\nu$ $\theta$$\lambda$$\tilde{\iota}$$\psi$$\iota$$\nu$} 
 \begin{itemize} 
\itemsep-.25em
\item{帖撒羅尼迦前書}\\ 2 我們為你們眾人常常感謝神，禱告的時候提到你們， 3 在神我們的父面前不住地記念你們因信心所做的工夫，因愛心所受的勞苦，因盼望我們主耶穌基督所存的忍耐。 4 被神所愛的弟兄啊，我知道你們是蒙揀選的， 5 因為我們的福音傳到你們那裡，不獨在乎言語，也在乎權能和聖靈並充足的信心。正如你們知道，我們在你們那裡，為你們的緣故是怎樣為人。 6 並且你們在\underline{大難 plural $\theta$$\lambda$$\acute{\iota}$$\psi$$\epsilon$$\iota$}之中蒙了聖靈所賜的喜樂，領受真道，就效法我們，也效法了主， 7 甚至你們做了馬其頓和亞該亞所有信主之人的榜樣。 8 因為主的道從你們那裡已經傳揚出來，你們向神的信心不但在馬其頓和亞該亞，就是在各處也都傳開了，所以不用我們說什麼話。
 \item{帖撒羅尼迦前書 3:1-10}\\
 1 我們既不能再忍，就願意獨自等在雅典， 2 打發我們的兄弟——在基督福音上做神執事的提摩太前去，堅固你們，並在你們所信的道上勸慰你們， 3 免得有人被諸般\underline{患難$\theta$$\lambda$$\acute{\iota}$$\psi$$\epsilon$$\sigma$$\iota$$\nu$}搖動。因為你們自己知道，我們受患難原是命定的。 4 我們在你們那裡的時候，預先告訴你們我們\underline{必受患難$\theta$$\lambda$$\acute{\iota}$$\beta$$\epsilon$$\sigma$$\theta$$\alpha$$\iota$$\tau$$o$ (to suffer tribulation)}，以後果然應驗了，你們也知道。 5 為此，我既不能再忍，就打發人去，要曉得你們的信心如何，恐怕那誘惑人的到底誘惑了你們，叫我們的勞苦歸於徒然。 6 但提摩太剛才從你們那裡回來，將你們信心和愛心的好消息報給我們，又說你們常常記念我們，切切地想見我們，如同我們想見你們一樣。 7 所以弟兄們，我們在一切困苦\underline{患難$\theta$$\lambda$$\acute{\iota}$$\psi$$\epsilon$$\iota$}之中，因著你們的信心就得了安慰。 8 你們若靠主站立得穩，我們就活了。 9 我們在神面前，因著你們甚是喜樂，為這一切喜樂，可用何等的感謝為你們報答神呢！ 10 我們晝夜切切地祈求，要見你們的面，補滿你們信心的不足。
 \item{帖後 1:3-10}\\ 3 弟兄們，我們該為你們常常感謝神，這本是合宜的，因你們的信心格外增長，並且你們眾人彼此相愛的心也都充足。 4 甚至我們在神的各教會裡為你們誇口，都因你們在所受的一切逼迫\underline{患難$\theta$$\lambda$$\acute{\iota}$$\psi$$\epsilon$$\sigma$$\iota$$\nu$}中，仍舊存忍耐和信心。 5 這正是神公義判斷的明證，叫你們可算配得神的國，你們就是為這國受苦。 6 神既是公義的，就必將患難報應那加\underline{患難$\theta$$\lambda$$\tilde{\iota}$$\psi$$\iota$$\nu$}給你們的人， 7 也必使你們這受患難的人與我們同得平安。那時，主耶穌同他有能力的天使從天上在火焰中顯現， 8 要報應那不認識神和那不聽從我主耶穌福音的人。 9 他們要受刑罰，就是永遠沉淪，離開主的面和他權能的榮光。 10 這正是主降臨，要在他聖徒的身上得榮耀，又在一切信的人身上顯為稀奇的那日子。我們對你們作的見證，你們也信了。
\end{itemize}
       

\subsubsection{希伯來書 10:32-39 $\theta$$\lambda$$\acute{\iota}$$\psi$$\epsilon$$\sigma$$\iota$$\nu$} 
 32 你們要追念往日蒙了光照以後，所忍受大爭戰的各樣苦難。 33 一面被毀謗，遭\underline{患難$\theta$$\lambda$$\acute{\iota}$$\psi$$\epsilon$$\sigma$$\iota$$\nu$}，成了戲景叫眾人觀看；一面陪伴那些受這樣苦難的人。 34 因為你們體恤了那些被捆鎖的人，並且你們的家業被人搶去，也甘心忍受，知道自己有更美、長存的家業。 35 所以，你們不可丟棄勇敢的心，存這樣的心必得大賞賜。 36 你們必須忍耐，使你們行完了神的旨意，就可以得著所應許的。 37 「因為還有一點點時候，那要來的就來，並不遲延。 38 只是義人 b必因信得生；他若退後，我心裡就不喜歡他。」 39 我們卻不是退後入沉淪的那等人，乃是有信心以致靈魂得救的人。
  
 \subsubsection{雅各書 1:26-27 $\theta$$\lambda$$\acute{\iota}$$\psi$$\epsilon$$\iota$}   
 26 若有人自以為虔誠，卻不勒住他的舌頭，反欺哄自己的心，這人的虔誠是虛的。 27 在神我們的父面前，那清潔沒有玷汙的虔誠就是看顧在\underline{患難$\theta$$\lambda$$\acute{\iota}$$\psi$$\epsilon$$\iota$}中的孤兒寡婦，並且保守自己不沾染世俗。   
 
 \subsubsection{啟示錄 1:9-11 $\theta$$\lambda$$\acute{\iota}$$\psi$$\epsilon$$\iota$, 2:8-11 $\theta$$\lambda$$\tilde{\iota}$$\psi$$\iota$$\nu$, 2:18-23 $\theta$$\lambda$$\tilde{\iota}$$\psi$$\iota$$\nu$, 7:13-17 $\theta$$\lambda$$\acute{\iota}$$\psi$$\epsilon$$\omega$$\zeta$} 
  \begin{itemize} 
\itemsep-.25em
\item{啟示錄 1:9-11}\\
 我約翰，就是你們的弟兄，和你們在耶穌的\underline{患難$\theta$$\lambda$$\acute{\iota}$$\psi$$\epsilon$$\iota$}、國度、忍耐裡一同有份，為神的道並為給耶穌作的見證，曾在那名叫拔摩的海島上。 10 當主日，我被聖靈感動，聽見在我後面有大聲音如吹號說： 11 「你所看見的當寫在書上，達於以弗所、士每拿、別迦摩、推雅推喇、撒狄、非拉鐵非、老底嘉那七個教會。」   
    
\item{啟示錄 2:8-11} \\ 
你要寫信給士每拿教會的使者說：『那首先的、末後的、死過又活的說： 9 我知道你的患難、你的貧窮，你卻是富足的；也知道那自稱是猶太人所說的毀謗話，其實他們不是猶太人，乃是撒旦一會的人。10 『你將要受的苦你不用怕。魔鬼要把你們中間幾個人下在監裡，叫你們被試煉，你們必受\underline{患難$\theta$$\lambda$$\tilde{\iota}$$\psi$$\iota$$\nu$}十日。你務要至死忠心，我就賜給你那生命的冠冕。 11 聖靈向眾教會所說的話，凡有耳的，就應當聽！得勝的，必不受第二次死的害。』

\item{啟示錄 2:18-23}   \\       
  18 「你要寫信給推雅推喇教會的使者說：『那眼目如火焰、腳像光明銅的神之子說： 19 我知道你的行為、愛心、信心、勤勞、忍耐，又知道你末後所行的善事比起初所行的更多。 20 然而，有一件事我要責備你，就是你容讓那自稱是先知的婦人耶洗別教導我的僕人，引誘他們行姦淫，吃祭偶像之物。 21 我曾給她悔改的機會，她卻不肯悔改她的淫行。 22 看哪，我要叫她病臥在床；那些與她行淫的人，若不悔改所行的，我也要叫他們同受\underline{大患難$\theta$$\lambda$$\tilde{\iota}$$\psi$$\iota$$\nu$}。 23 我又要殺死她的黨類 a，叫眾教會知道我是那察看人肺腑心腸的，並要照你們的行為報應你們各人。   
 \item{啟示錄 7:13-17} \\    
  13 長老中有一位問我說：「這些穿白衣的是誰？是從哪裡來的？」 14 我對他說：「我主，你知道。」他向我說：「這些人是從\underline{大患難$\theta$$\lambda$$\acute{\iota}$$\psi$$\epsilon$$\omega$$\zeta$}中出來的，曾用羔羊的血把衣裳洗白淨了。 15 所以，他們在神寶座前，晝夜在他殿中侍奉他；坐寶座的要用帳幕覆庇他們。 16 他們不再飢、不再渴，日頭和炎熱也必不傷害他們， 17 因為寶座中的羔羊必牧養他們，領他們到生命水的泉源；神也必擦去他們一切的眼淚。」   
\end{itemize}


\subsection{神的印}


\subsection{兽的印记}


\subsection{天上的星，海边的沙}

創世記 15:5
於是領他走到外邊，說：「你向天觀看，數算眾星，能數得過來嗎？」又對他說：「你的後裔將要如此。」

創世記 22:17 論福，我必賜大福給你；論子孫，我必叫你的子孫多起來，如同天上的星，海邊的沙；你子孫必得著仇敵的城門；

創世記 26:4
我要加增你的後裔，像天上的星那樣多，又要將這些地都賜給你的後裔，並且地上萬國必因你的後裔得福。

創世記 28:14
你的後裔必像地上的塵沙那樣多，必向東西南北開展。地上萬族必因你和你的後裔得福。

創世記 32:12
你曾說：『我必定厚待你，使你的後裔如同海邊的沙，多得不可勝數。』」

出埃及記 32:13
求你記念你的僕人亞伯拉罕、以撒、以色列，你曾指著自己起誓說：『我必使你們的後裔像天上的星那樣多，並且我所應許的這全地，必給你們的後裔，他們要永遠承受為業。』」

申命記 1:10
耶和華你們的神使你們多起來，看哪，你們今日像天上的星那樣多！
申命記 10:22
你的列祖七十人下埃及，現在耶和華你的神使你如同天上的星那樣多。

撒母耳記下 17:11 依我之計，不如將以色列眾人從但直到別是巴，如同海邊的沙那樣多聚集到你這裡來，你也親自率領他們出戰。

列王紀上 4:20
猶大人和以色列人如同海邊的沙那樣多，都吃喝快樂。
列王紀上 4:29
神賜給所羅門極大的智慧、聰明和廣大的心，如同海沙不可測量

歷代志上 27:23
以色列人二十歲以內的，大衛沒有記其數目，因耶和華曾應許說必加增以色列人如天上的星那樣多。
歷代志下 1:9
耶和華神啊，現在求你成就向我父大衛所應許的話，因你立我做這民的王，他們如同地上塵沙那樣多。

尼希米記 9:23 你也使他們的子孫多如天上的星，帶他們到你所應許他們列祖進入得為業之地。

以賽亞書 48:19 你的後裔也必多如海沙，你腹中所生的也必多如沙粒，他的名在我面前必不剪除，也不滅絕。」

耶利米書 33:22
天上的萬象不能數算，海邊的塵沙也不能斗量，我必照樣使我僕人大衛的後裔和侍奉我的利未人多起來。」

何西阿書 1:10
「然而，以色列的人數必如海沙，不可量，不可數。從前在什麼地方對他們說『你們不是我的子民』，將來在那裡必對他們說『你們是永生神的兒子』。


羅馬書 9:27
以賽亞指著以色列人喊著說：「以色列人雖多如海沙，得救的不過是剩下的餘數

希伯來書 11:12
所以從一個彷彿已死的人就生出子孫，如同天上的星那樣眾多，海邊的沙那樣無數。

啟示錄 12:17
龍向婦人發怒，去與她其餘的兒女爭戰，這兒女就是那守神誡命、為耶穌作見證的。那時龍就站在海邊的沙上。




\subsection{地開口}
創世記 4:9-12
耶和華對該隱說：「你兄弟亞伯在哪裡？」他說：「我不知道。我豈是看守我兄弟的嗎？」 10 耶和華說：「你做了什麼事呢？你兄弟的血有聲音從地裡向我哀告。 11 地開了口，從你手裡接受你兄弟的血，現在你必從這地受咒詛。 12 你種地，地不再給你效力；你必流離飄蕩在地上。」

申命記 11:6
也沒有看見他怎樣待魯本子孫以利押的兒子大坍、亞比蘭，地怎樣在以色列人中間開口，吞了他們和他們的家眷，並帳篷與跟他們的一切活物

民數記 16:29-32
29這些人死若與世人無異，或是他們所遭的與世人相同，就不是耶和華打發我來的。 30倘若耶和華創作一件新事，使地開口，把他們和一切屬他們的都吞下去，叫他們活活地墜落陰間，你們就明白這些人是藐視耶和華了。」摩西剛說完了這一切話，他們腳下的地就開了口，把他們和他們的家眷，並一切屬可拉的人丁、財物，都吞下去。

民數記 26:10
以利押的眾子是尼母利、大坍、亞比蘭。這大坍、亞比蘭，就是從會中選召的，與可拉一黨同向耶和華爭鬧的時候也向摩西、亞倫爭鬧。 10地便開口吞了他們，和可拉、可拉的黨類一同死亡。那時火燒滅了二百五十個人，他們就做了警戒。 11然而可拉的眾子沒有死亡

啟示錄 12:16
15蛇就在婦人身後，從口中吐出水來，像河一樣，要將婦人沖去。 16地卻幫助婦人，開口吞了從龍口吐出來的水（河水）。
  
  
  
  
  
  
  
  
  
  
  
  
  
  
  
  
  
  